% Long Note: Original attachment comments are preserved verbatim below, including resolved and superseded requests.
% Long Note: Original line anchors identify the supplied attachment; these comments do not replace the current manuscript.
% Long Note: Dispositions and new editorial responses remain beside the relevant current text and in the comment-response audit.

% Long Note: Original attachment line 3.
% Prevent pdflatex object-stream overflow with complex vector figures.
% Long Note: Done [T01]. Original compression workaround retained with an engine guard;
% native multiproject compilation is checked separately.

% Long Note: Original attachment line 27.
% Uncomment if line numbers are required.

% Long Note: Original attachment line 28.
% \usepackage[mathlines]{lineno}\linenumbers
% Long Note: Done [T02]. Optional line-numbering switch retained, not activated without a
% submission requirement.

% Long Note: Original attachment line 44.
% \titleformat{\section}{\large\bfseries}{\thesection.}{0.6em}{}

% Long Note: Original attachment line 45.
% \titleformat{\subsection}{\normalsize\bfseries}{\thesubsection.}{0.6em}{}
% Long Note: Done [T03]. Commented formatting alternatives retained; active headings are
% unnumbered.

% Long Note: Original attachment line 58.
%
% Long Note: Done [T04]. Macro whitespace-control comment retained.

% Long Note: Original attachment line 71.
% Nature Machine Intelligence
% Long Note: Done [T05]. Venue/style note retained; no acceptance claim is inferred.

% Long Note: Original attachment line 75.
% Pangenome Foundation Model capture
% Long Note: Done [T06]. Title scopes the contribution to topology-native representation
% learning.

% Long Note: Original attachment line 76.
% Need to add HGSVC and HPRC to the author list

% Long Note: Original attachment line 82.
% HGSVC, HPRC,
% Long Note: Author input needed [A01]. Exact approved consortium author wording, membership
% and approval are unavailable.

% Long Note: Original attachment line 101.
% NEW STRUCTURES

% Long Note: Original attachment line 103.
% Title

% Long Note: Original attachment line 104.
% Abstract

% Long Note: Original attachment line 106.
% Introduction

% Long Note: Original attachment line 108.
% Results

% Long Note: Original attachment line 109.
%   1. Graph-native representations add complementary information to regulatory-element prediction

% Long Note: Original attachment line 110.
%   2. Graph-native representations contribute strongly to structural-variant breakpoint classification

% Long Note: Original attachment line 111.
%   3. Frozen graph representations generalize across chromosomes and transfer across pangenome resources

% Long Note: Original attachment line 112.
%   4. Neighborhood context distinguishes learned graph representations from degree-based shortcuts

% Long Note: Original attachment line 114.
% Discussion

% Long Note: Original attachment line 116.
% Methods

% Long Note: Original attachment line 117.
%   Pangenome resources and graph preparation

% Long Note: Original attachment line 118.
%   Masked graph-connection pretraining

% Long Note: Original attachment line 119.
%   PangenomeFM architecture

% Long Note: Original attachment line 120.
%   Strict and one-hop graph context

% Long Note: Original attachment line 121.
%   Chromosome-held-out evaluation

% Long Note: Original attachment line 122.
%   Cross-resource and cross-release transfer

% Long Note: Original attachment line 123.
%   Sequence embeddings and downstream classifiers

% Long Note: Original attachment line 124.
%   Baselines, ablations and capacity

% Long Note: Original attachment line 125.
%   Statistical analysis

% Long Note: Original attachment line 126.
%   Exploratory path/QTL/GWAS analyses

% Long Note: Original attachment line 128.
% Data availability

% Long Note: Original attachment line 129.
% Code availability

% Long Note: Original attachment line 130.
% Acknowledgements

% Long Note: Original attachment line 131.
% Funding

% Long Note: Original attachment line 132.
% Author contributions

% Long Note: Original attachment line 133.
% References

% Long Note: Original attachment line 135.
% Supplementary Information
% Long Note: Done [T07]. Historical outline retained as comments; the revised method-led
% structure supersedes its ordering.

% Long Note: Original attachment line 146.
% Structure: biological representation problem --> graph-native idea --> cCRE --> SV --> generalization --> control --> scoped conclusion
% Long Note: Done [T08]. Abstract revised around the representation problem, design, evidence
% and attribution limit.

% Long Note: Original attachment line 153.
% Long and Tomoya, here are what I think about the layout of the introduction: representation of DNA sequences and human genomes (linear->graph - capturing diversity, complex variations and regions), existing DNA foundation models and their applications (this can be swapped to the beginning of this session), the need of non-linear/graph representation of genomes and foundation models  to capture information in graph genomes or pangnomes, previous efforts of building pangenome foundation models and what they lack of, -> introduce pangenomeFM, improvement over previous work, applications of using pangenomeFM for downstream tasks.
% Long Note: Done [S01]. Introduction progression implemented; R01.
% Long Note: Done [R01]. Introduction follows the six-paragraph representation-to-method
% progression; see sec:intro.

% Long Note: Original attachment line 157.
% Long, the following paper should be cited (add other related papers if any) to stress the importance of using graph genomes for functional genomics analysis

% Long Note: Original attachment line 158.
% https://www.nature.com/articles/s41467-026-73663-3
% Long Note: Done [S02]. Requested functional-genomics citation verified; R02.
% Long Note: Done [R02]. Macias-Velasco et al. is cited with DOI 10.1038/s41467-026-73663-3;
% its genome/pipeline scope is stated accurately.

% Long Note: Original attachment line 161.
% [LONG NOTE]

% Long Note: Original attachment line 162.
% Para 1: linear to graph: what linear has, what it loses, and why graph could enable the diversity

% Long Note: Original attachment line 163.
% Para 2: why the representation matters in biology

% Long Note: Original attachment line 164.
% para 3: talks about the DNA foundation models (DNABERT,etc.)

% Long Note: Original attachment line 165.
% para 4: talks about pangenome topology features which sequence tokens do not present/have

% Long Note: Original attachment line 166.
% para 5: previous work & gaps (DeepGene, PangenomeX, etc.)

% Long Note: Original attachment line 167.
% para 6: introduce PangenomeFM + evidence/result (main generalization, etc.)
% Long Note: Done [S03]. Six-paragraph outline implemented in sec:intro.

% Long Note: Original attachment line 172.
% Intro should start with the representation of human genomes, genetic variation should be mentioned later
% Long Note: Done [S04]. Representation-first opening retained; R03.
% Long Note: Done [R03]. The Introduction opens with human-genome representation and
% coordinate systems before discussing variation.

% Long Note: Original attachment line 183.
% Long, the layout looks more like a CS paper, we need to change it to the format of Nature Machine Intelligence or Nature Methods (check some recent papers in those journals and follow their layout)
% Long Note: Done [S05]. NMI examples reviewed and structure revised; R04.
% Long Note: Done [R04]. All four supplied NMI examples informed the method-led structure;
% four main figures and two main tables retain key controls.

% Long Note: Original attachment line 187.
% Long, this figure should be cited in Introduction, where you can merge the current related work and background into Introduction
% Long Note: Done [S06]. Paradigm figure cited in integrated Introduction; R05.
% Long Note: Done [R05]. The Introduction integrates background/prior work and cites
% fig:paradigms; no separate Related Work section is needed.

% Long Note: Original attachment line 190.
% TODO(figure): redesign with keyword-only labels. End each paradigm at

% Long Note: Original attachment line 191.
% its learned output; explain downstream reuse in the caption rather than a

% Long Note: Original attachment line 192.
% PangenomeFM-only output column.
% Long Note: Done [T09]. Parallel output figure redesigned and verified; R06.
% Long Note: Done [R06]. The redesigned keyword-led paradigm figure ends every row at its
% learned representation; reuse is explained in the caption.

% Long Note: Original attachment line 264.
% Long, add results on predicting cCREs and ENCODE data?

% Long Note: Original attachment line 265.
% add citation of this paper 

% Long Note: Original attachment line 266.
% https://www.nature.com/articles/s41467-026-73663-3
% Long Note: Done [S07]. cCRE/ENCODE result and requested citation retained; R07/R02.

% Long Note: Original attachment line 268.
% cCRE is now the first standalone Result. The existing binary cCRE

% Long Note: Original attachment line 269.
% result is included below. TODO(new analysis): if ENCODE class labels are cleanly

% Long Note: Original attachment line 270.
% recoverable, add PLS/pELS/dELS/CTCF-only stratification and/or topology gain by

% Long Note: Original attachment line 271.
% local graph/variant complexity. Do not assume which class benefits most.
% Long Note: Done [R07]. Binary cCRE, PLS/pELS/dELS/CTCF-only and fixed complexity strata are
% reported in sec:ccre-result. The old first-Result placement is superseded by the
% method-first order.

% Long Note: Original attachment line 282.
% TODO(new analysis): add cCRE class / complexity stratification here if completed.

% Long Note: Original attachment line 283.
% Recommended tests: PLS, pELS, dELS, CTCF-only; and/or low/medium/high local

% Long Note: Original attachment line 284.
% graph complexity using thresholds fixed independently of model performance.
% Long Note: Done [S08]. Four subclasses and fixed complexity strata reported; R07.

% Long Note: Original attachment line 285.
% Optional biological vignette: one transparently selected locus with graph,

% Long Note: Original attachment line 286.
% annotation track, and predictions; avoid selecting solely by maximal gain.
% Long Note: Optional, not done [S09]. No selected-locus vignette is claimed; it needs
% verified tracks, predictions and a predeclared selection rule.

% Long Note: Original attachment line 320.
% Keep this table:

% Long Note: Original attachment line 321.
% For final NMI submission, consider moving the table to Supplementary once the

% Long Note: Original attachment line 322.
% standalone biological figures carry these values clearly.
% Long Note: Done [T10]. Primary numerical table preserved in Supplement; main biological
% figures retain its outcomes.

% Long Note: Original attachment line 343.
% Is it possible to improve SV genotyping (especially in complex and copy number variant regions)?
% Long Note: Evidence open [M01]. Genotyping assessment answered, but improved calls are not
% established; R08.
% Long Note: Evidence open [R08]. The genotyping question is answered by the completed COSIGT
% assessment, but improved genotype calls, especially complex/CNV calls, are NOT
% demonstrated. Quality prediction is not genotype-concordance improvement; see
% sec:generalization-result.

% Long Note: Original attachment line 373.
% Long, this subsection title is too cs/graph-heavy, we need to rephrase
% Long Note: Done [S10]. Heading revised without unsupported learned-neighborhood inference;
% R09.
% Long Note: Done [R09]. The heading now reads Structural information and pretraining make
% different contributions; no neighborhood-learning superiority is assumed.

% Long Note: Original attachment line 381.
% TODO: because the degree-baseline values differ between strict and one-hop

% Long Note: Original attachment line 382.
% analyses, state explicitly whether node degree was recomputed after context expansion.

% Long Note: Original attachment line 383.
% HIGH PRIORITY CONTROL: recompute degree sum / preferential attachment on the

% Long Note: Original attachment line 384.
% actual expanded one-hop graph visible to the baseline. Do not strengthen the

% Long Note: Original attachment line 385.
% "neighborhood organization" claim until this control is verified.

% Long Note: Original attachment line 410.
% Long, I feel that the first (or the first two) subsessions should be dedicated to novel components of PangenomeFM (and highlight those key components)
% Long Note: Done [S11]. Method components lead Results and Methods; R10.
% Long Note: Done [R10]. The first Results section explains the design; Methods begins with
% representation/coupled encoder, objective/optimization and frozen reuse.
% Established components are attributed.

% Long Note: Original attachment line 416.
% Long, specify how many graph genomes, etc.
% Long Note: Partial, evidence open [S12]. Graph/cache and path-resource counts specified;
% exact donor provenance of the SV-resolution training graph remains unverified.

% Long Note: Original attachment line 431.
% TODO(authors): state the connected:unconnected sampling ratio and the resulting

% Long Note: Original attachment line 432.
% chance-level AUPRC. This is needed to interpret values such as 0.5997.

% Long Note: Original attachment line 448.
% TODO: specify the implemented gate function, pair-scoring architecture,

% Long Note: Original attachment line 449.
% and the exact focal/class-adaptive form of w_{uv}; these are required for reproducibility.

% Long Note: Original attachment line 450.
% Suggested final wording must be copied from the actual archived code/config, not inferred.
% Long Note: Done [M04]. Shared-state gate, ordered scorer and exact adopted focal loss
% specified.

% Long Note: Original attachment line 458.
% We may need to talk about loss functions and optimization in PangenomFM
% Long Note: Done [M05]. Loss and optimization specified; R11.
% Long Note: Done [R11]. sec:objective specifies the ordered scorer, exact focal loss and
% class weighting, AdamW, schedule, clipping and stopping criteria.

% Long Note: Original attachment line 480.
% TODO(authors): once the class ratio is added above, report the corresponding

% Long Note: Original attachment line 481.
% chance AUPRC here and add it as a reference line to the generalization figure.
% Long Note: Done [M03]. Prevalence/chance levels and reference line reported; actual
% context-specific ratios distinguished.

% Long Note: Original attachment line 487.
% TODO(authors): state exactly which graph is used to compute degree under each

% Long Note: Original attachment line 488.
% context. Recompute degree on the expanded one-hop graph if this was not done in

% Long Note: Original attachment line 489.
% the original baseline. The main-text interpretation must be updated if the

% Long Note: Original attachment line 490.
% strict-to-one-hop reversal changes after this fair control.
% Long Note: Done [M02]. Visible-graph baseline replay verified; historical masking
% limitations remain explicit.

% Long Note: Original attachment line 493.
% TODO(authors): give the numeric cutoffs used to define the three complexity strata

% Long Note: Original attachment line 494.
% and the number of regions/examples in each stratum.

% Long Note: Original attachment line 497.
% TODO(authors): the principal configuration (48/2) is not part of this sweep, and the

% Long Note: Original attachment line 498.
% 96/4 setting reports higher strict AUPRC than the principal HPRC R2 model. Preferably

% Long Note: Original attachment line 499.
% run 48/2 under the exact diagnostic sweep protocol; otherwise state explicitly why

% Long Note: Original attachment line 500.
% the sweep and the main runs are not directly comparable.

% Long Note: Original attachment line 508.
% Tomoya and Long, These results should be included as a separate subsection in Results, including detailed interpretations of these results,

% Long Note: Original attachment line 509.
% RECONCILED IN THIS REVISION: the scientifically relevant coverage/truncation result

% Long Note: Original attachment line 510.
% is summarized in the first biological Results section; detailed QC remains here

% Long Note: Original attachment line 511.
% and in Supplementary so it does not consume one of the <=5 headline Results sections.
% Long Note: Done [S13]. Dedicated coverage/truncation Results heading and interpretation;
% R12.
% Long Note: Done [R12]. A dedicated sequence-coverage Results subsubsection reports full
% segment coverage and 26.3 percent end sampling, with detailed QC in the
% supplement.

% Long Note: Original attachment line 516.
% TODO: name the downstream classifier family and its hyperparameters,

% Long Note: Original attachment line 517.
% regularization, model-selection procedure, and whether any threshold tuning used

% Long Note: Original attachment line 518.
% validation chromosomes only.
% Long Note: Done [M08]. Original and later probe settings and validation-only selection
% distinguished.

% Long Note: Original attachment line 520.
% TODO: clarify whether 60,685 applies to each split or to validation and

% Long Note: Original attachment line 521.
% test combined. Do not leave this ambiguous in the submission version.
% Long Note: Done [M09]. cCRE validation and test counts explicitly apply to each split.

% Long Note: Original attachment line 526.
% TODO: clarify whether 34,781 applies to each split or to validation and

% Long Note: Original attachment line 527.
% test combined.
% Long Note: Done [M10]. INS/DEL validation and test counts explicitly apply to each split.

% Long Note: Original attachment line 552.
% TODO: state the actual experimental reason the QTL overlap analysis was

% Long Note: Original attachment line 553.
% restricted to chr8 (for example, source-data availability or pilot scope). Do NOT

% Long Note: Original attachment line 554.
% use an unrelated chr8 literature example as a post-hoc justification.
% Long Note: Done [M11]. Actual chromosome-8 pilot scope stated without a post-hoc biological
% rationale.

% Long Note: Original attachment line 567.
% Long, please prepare a github under shilab repo for your code, should we share the pretrained model there or huggingface etc.?
% Long Note: Author input needed [A02]. Current repository disclosed; Shi-lab destination,
% release license and checkpoint hosting remain unconfirmed.

% Long Note: Original attachment line 574.
% acknowledge NHGRI U24 and NIGMS R01 for funding
% Long Note: Author input needed [A03]. Exact grant award numbers, recipients and required
% acknowledgement wording remain unconfirmed.

% Long Note: Original attachment line 588.
% Long, Supplementary information should read like a manuscript, including its own sessions, text, figures, tables, and references
% Long Note: Done [S14]. Supplement has its own narrative and references; R13.
% Long Note: Done [R13]. The supplement has its own narrative sections, figures, tables and
% S-numbered references; in-text S labels now match its bibliography.

% Long Note: Original attachment line 689.
% TODO(authors): give exact graph-complexity thresholds and counts here once recovered.
% Long Note: Done [M06]. Six-component complexity definition, fixed cutoffs and region counts
% supplied.

% Long Note: Original attachment line 694.
% TODO(authors): the principal 48/2 configuration is not included in this sweep.

% Long Note: Original attachment line 695.
% Preferably add the exact 48/2 diagnostic run; otherwise explain why these rows

% Long Note: Original attachment line 696.
% are not directly comparable with the principal HPRC R2 experiment.
% Long Note: Done [M07]. Archived matched 48/2 evidence recovered; fold-pooled versus
% window-mean distinction stated.
